%----------------------------------------------------------------------------------------
%    PACKAGES AND THEMES
%----------------------------------------------------------------------------------------

\documentclass[aspectratio=169,xcolor=dvipsnames]{beamer}
\usetheme{SimplePlus}
\usepackage{subcaption}
\usepackage{hyperref}
\usepackage{amsmath}
\usepackage{xcolor}
\usepackage{graphicx}
\usepackage{booktabs}
\usepackage{tikz}

%----------------------------------------------------------------------------------------
%    TITLE PAGE
%----------------------------------------------------------------------------------------

\title{SSM-CGM: Interpretable State-Space Forecasting Model of Continuous Glucose Monitoring for Personalized Diabetes Management}
\subtitle{Literature Review and Analysis}

\author{DAIT DEHANE Yacine}

\institute
{
    Supervised by \\
    - Pr Henni Khadidja (TÉLUQ)\\
    - Dr Khalid Benabdeslem (Université de Lyon)
}
\date{\today}

%----------------------------------------------------------------------------------------
%    PRESENTATION SLIDES
%----------------------------------------------------------------------------------------

\begin{document}

% ----------------------------------------------------
% Frame 1: Title Slide
% ----------------------------------------------------
\begin{frame}
    \titlepage
\end{frame}
% ----------------------------------------------------
% Frame 2: Presentation Roadmap
% ----------------------------------------------------
\begin{frame}{Presentation Roadmap}
 
    
     
    \begin{columns}[t]
        \begin{column}{0.48\textwidth}
            \begin{block}{\textbf{Part I: The Challenge}}
                \begin{itemize}
                    \item Global diabetes crisis
                    \item CGM data opportunities
                    \item The interpretability gap
                    \item AI-READI dataset
                \end{itemize}
            \end{block}
            
            \vspace{0.3em}
            
            \begin{block}{\textbf{Part II: Baseline Approach}}
                \begin{itemize}
                    \item Temporal Fusion Transformer
                    \item Strengths and limitations
                    \item Computational bottleneck
                \end{itemize}
            \end{block}
        \end{column}
        
        \begin{column}{0.48\textwidth}
            \begin{block}{\textbf{Part III: SSM-CGM Solution}}
                \begin{itemize}
                    \item Mamba-based architecture
                    \item Linear-time complexity
                    \item State-space mechanics
                    \item Preserved interpretability
                \end{itemize}
            \end{block}
            
            \vspace{0.3em}
            
            \begin{block}{\textbf{Part IV: Validation \& Impact}}
                \begin{itemize}
                    \item Experimental results
                    \item Interpretability analysis
                    \item Counterfactual forecasting
                    \item Future directions
                \end{itemize}
            \end{block}
        \end{column}
    \end{columns}
    
    \vspace{0.5em}
    
  
\end{frame}

% ----------------------------------------------------
% Frame 2: Introduction - The Diabetes Crisis
% ----------------------------------------------------
\begin{frame}{Introduction: The Global Diabetes Challenge}
    \begin{block}{The Magnitude of the Problem}
        \begin{itemize}
            \item \textbf{Global Impact:} 580+ million adults worldwide live with diabetes
            \item \textbf{Complications:} Poor glucose control leads to cardiovascular disease, neuropathy, retinopathy, kidney failure
            \item \textbf{Management Complexity:} Glucose fluctuates based on meals, activity, stress, medications, circadian rhythms
        \end{itemize}
    \end{block}
    
    \begin{block}{Continuous Glucose Monitoring: A Game Changer}
        \begin{itemize}
            \item \textbf{Technology:} CGM devices (Dexcom G6) measure glucose every 5--15 minutes
            \item \textbf{Duration:} Continuous monitoring up to 14 days
            \item \textbf{Clinical Metric:} Time-in-Range (TIR: 70--180 mg/dL) correlates with reduced complications
            \item \textbf{Data Volume:} Thousands of data points per patient
        \end{itemize}
    \end{block}
\end{frame}

% ----------------------------------------------------
% Frame 3: The Core Problem
% ----------------------------------------------------
\begin{frame}{The Challenge: From Data to Action}
    \begin{columns}
        \begin{column}{0.5\textwidth}
            \begin{block}{What We Have}
                \begin{itemize}
                    \item Dense, high-resolution CGM data streams
                    \item Wearable activity data (heart rate, respiration, steps, sleep)
                    \item Clinical information (age, diabetes status, medications)
                \end{itemize}
            \end{block}
            
            \begin{block}{What We Need}
                \begin{itemize}
                    \item Accurate near-future glucose predictions 
                    \item \textbf{Interpretable} insights into what drives glucose changes
                    \item \textbf{Actionable} Patient-specific insights
                \end{itemize}
            \end{block}
        \end{column}
        
        \begin{column}{0.5\textwidth}
            \begin{block}{The Opportunity(The gap)}
                \textbf{Existing forecasting models:}
                \begin{itemize}
                    \item Achieve good predictive accuracy
                    \item But operate as \textbf{black boxes}
                    \item Lack clinical interpretability
                    \item Cannot explain \textit{why} predictions are made
                    \item Don't support "what-if" scenario planning
                \end{itemize}
            \end{block}
            
             
            \begin{alertblock}{Research Goal}
                Develop a model that is \textbf{accurate}, \textbf{interpretable}, and enables \textbf{personalized} behavioral guidance for diabetes management.
            \end{alertblock}
        \end{column}
    \end{columns}
\end{frame}

% ----------------------------------------------------
% TRANSITION 1: From Problem to Data
% ----------------------------------------------------
\begin{frame}[plain]
    \centering
    \vfill
    \begin{beamercolorbox}[center,shadow=true,rounded=true]{title}
        \usebeamerfont{title}
        \textbf{Part I: Understanding the Challenge}\\
        \vspace{0.5em}
        \large Problem $\to$ Available Data $\to$ Previous Solutions
    \end{beamercolorbox}
    \vfill
\end{frame}

% ----------------------------------------------------
% Frame 4: Dataset MOVED HERE
% ----------------------------------------------------
\begin{frame}{Dataset: AI-READI Cohort}
    \begin{columns}
        \begin{column}{0.5\textwidth}
            \begin{block}{AI-READI: \textsc{\textit{Artificial Intelligence Ready and Equitable Atlas for Diabetes Insights}}}
                \textbf{Cohort Characteristics:}
                \begin{itemize}
                    \item \textbf{Size:} 741 participants
                    \item \textbf{Age Range:} 40+ years (mean: 60 $\pm$ 11 years)
                    \item \textbf{Glycemic Status:}
                    \begin{itemize}
                        \item Healthy: 275 (37\%)
                        \item Prediabetes: 163 (22\%)
                        \item T2D (oral meds): 213 (29\%)
                        \item T2D (insulin): 90 (12\%)
                    \end{itemize}
                    \item \textbf{Duration:} 8--10 days per participant
                    \item \textbf{Total Measurements:} +2 million data points
                \end{itemize}
            \end{block}
        \end{column}
        
        \begin{column}{0.5\textwidth}
            \begin{block}{Multimodal Data Sources}
                \textbf{CGM Device:} Dexcom G6 (5-min sampling)
                
                \vspace{0.05em}
                \textbf{Wearable Device:} Garmin Vivosmart 5
                \begin{itemize}
                    \item Heart rate
                    \item Respiration rate
                    \item Oxygen saturation (SpO$_2$)
                    \item Steps
                    \item Stress level
                    \item Calories burned
                    \item Sleep stage
                \end{itemize}
                
                \vspace{0.05em}
                \textbf{Clinical Lab Values:}
                \begin{itemize}
                    \item HbA1c (mean: 6.0 $\pm$ 1.3\%)
                    \item Lipid panels, troponin, insulin
                \end{itemize}
            \end{block}
        \end{column}
    \end{columns}
\end{frame}

\begin{frame}
    \vspace{0.5em}
    \begin{block}{Meal Detection Challenge \& Solution}
        \textbf{Problem:} AI-READI lacks meal annotations \\
        \textbf{Solution:} Trained separate meal-detection model on CGMacros dataset (45 subjects with labeled meals)
        \begin{itemize}
            \item Dilated CNN + Bi-LSTM architecture with trapezoidal label smoothing
            \item \textbf{Performance:} 80\% Recall, 72\% Precision (sufficient for meal flag proxy)
            \item Deployed to generate predicted meal flags for AI-READI
        \end{itemize}
    \end{block}
\end{frame}

% ----------------------------------------------------
% TRANSITION 2: From Data to Baseline Model
% ----------------------------------------------------
\begin{frame}[plain]
    \centering
    \vfill
    \begin{beamercolorbox}[center,shadow=true,rounded=true]{title}
        \usebeamerfont{title}
        \textbf{Part II: The Baseline Approach}\\
        \vspace{0.5em}
        \large What exists? $\to$ What are its limitations?
    \end{beamercolorbox}
    \vfill
    
    \vspace{1em}
    \normalsize
    \textit{Now that we understand the problem and available data,\\
    let's examine how current models tackle CGM forecasting...}
\end{frame}

% ----------------------------------------------------
% Frame 5: Baseline Model - TFT (Conceptual Overview)
% ----------------------------------------------------
\begin{frame}{Baseline Model: Temporal Fusion Transformer (TFT)}
    \begin{block}{TFT Architecture Overview}
        The Temporal Fusion Transformer combines:
        \begin{itemize}
            \item \textbf{Variable Selection Networks:} Learn which features matter
            \item \textbf{LSTM layers:} Process sequential data for local patterns
            \item \textbf{Multi-head Self-Attention:} Capture global dependencies
            \item \textbf{Quantile prediction:} Output prediction intervals with uncertainty
        \end{itemize}
    \end{block}
    
 
\begin{center}
\Large
$\hat{y}_i(q, t, \tau) 
    = f_q\Big( \tau,\;
     {y_{i,t-k:t}},\;
     {z_{i,t-k:t}},\;
     {x_{i,t:t+\tau}},\;
      {s_i}
    \Big)$
\end{center}
\begin{center}

\footnotesize
Formula from Lim et al., Temporal Fusion Transformers (2021)~\cite{lim2021temporal}.
\end{center}
\end{frame}

% ----------------------------------------------------
% Frame 6: TFT Limitations (MOTIVATION)
% ----------------------------------------------------
\begin{frame}{Why TFT Falls Short for Long-Context CGM}
    \begin{block}{The Challenge: Long Sequences}
        CGM requires processing \textbf{long sequences} (12--48 hours = 144--576 time steps). TFT faces computational bottlenecks:
        
        \vspace{0.3em}
        \begin{columns}
            \begin{column}{0.48\textwidth}
                \textbf{LSTM Limitations:}
                \begin{itemize}
                    \item Sequential (cannot parallelize)
                    \item Vanishing gradients
                    \item Struggles with long dependencies
                \end{itemize}
            \end{column}
            
            \begin{column}{0.48\textwidth}
                \textbf{Attention Bottleneck:}
                \begin{itemize}
                    \item Quadratic complexity: $\mathcal{O}(L^2)$
                    \item 48h = 576 steps $\rightarrow$ 331k operations
                    \item High memory, slow inference
                \end{itemize}
            \end{column}
        \end{columns}
    \end{block}
    
    \vspace{0.5em}
    \begin{alertblock}{The Need for a Better Solution}
        We need a model that:
        \begin{itemize}
            \item Handles long contexts efficiently ($\mathcal{O}(L)$ instead of $\mathcal{O}(L^2)$)
            \item Captures long-range dependencies effectively
            \item Maintains interpretability for clinical use
        \end{itemize}
    \end{alertblock}
\end{frame}

% ----------------------------------------------------
% TRANSITION 3: From Limitations to Solution
% ----------------------------------------------------
\begin{frame}[plain]
    \centering
    \vfill
    \begin{beamercolorbox}[center,shadow=true,rounded=true]{title}
        \usebeamerfont{title}
        \textbf{Part III: The Proposed Solution}\\
        \vspace{0.5em}
        \large SSM-CGM: Keeping interpretability, gaining efficiency
    \end{beamercolorbox}
    \vfill
    
    \vspace{1em}
    \normalsize
    \begin{itemize}
        \item TFT showed us \textbf{what works}: interpretable architecture, quantile forecasting
        \item TFT showed us \textbf{what doesn't}: slow, quadratic complexity for long sequences
        \item SSM-CGM: \textbf{Keep the good, fix the bad}
    \end{itemize}
    
    \vspace{1em}
    \textit{Let's see how Mamba-based state-space models solve the bottleneck...}
\end{frame}

% ----------------------------------------------------
% Frame 7: SSM-CGM Solution (Conceptual)
% ----------------------------------------------------
\begin{frame}{SSM-CGM: The Proposed Solution}
    \begin{block}{Core Innovation}
        SSM-CGM retains TFT's interpretable interface (VSNs, multi-horizon head) but replaces the inefficient temporal processor with \textbf{Mamba}, a selective State-Space Model.
    \end{block}
    
    \begin{center}
        \begin{tabular}{p{0.25\textwidth} p{0.3\textwidth} p{0.3\textwidth}}
            \toprule
            \textbf{Component} & \textbf{TFT} & \textbf{SSM-CGM} \\
            \midrule
            Variable Selection & VSN & VSN (kept) \\
            \midrule
            Local Processing & LSTM (sequential) & \textbf{Stacked Mamba} (parallel) \\
            \midrule
            Global Context & Self-Attention ($\mathcal{O}(L^2)$) & \textbf{Light Mamba} ($\mathcal{O}(L)$) \\
            \midrule
            Prediction Head & Multi-horizon quantile & Multi-horizon quantile \\
            \bottomrule
        \end{tabular}
    \end{center}
\end{frame}

\begin{frame}{}
    \vspace{1em}
    \begin{block}{Mamba's Key Advantages}
        \begin{itemize}
            \item \textbf{Linear-time complexity:} $\mathcal{O}(L)$ instead of $\mathcal{O}(L^2)$ -- dramatically faster for long sequences
            \item \textbf{Selective state-space:} Dynamically learns what to remember/forget at each step
            \item \textbf{Efficient memory:} Small footprint enables real-time, on-device deployment
            \item \textbf{Parallelizable:} Unlike LSTMs, can process sequences in parallel during training
        \end{itemize}
    \end{block}
    
   \begin{center}
\Large
$\hat{y}_{t+1:t+H} 
= f_\theta \Big( 
    \textcolor{ForestGreen}{y_{t-L+1:t}},\;
    \textcolor{violet}{x^{(hist)}_{t-L+1:t}},\;
    \textcolor{BrickRed}{x^{(fut)}_{t+1:t+H}},\;
    \textcolor{Orange}{s}
\Big)$
\end{center}

\begin{center}
\footnotesize
Formula from Shakson et al.~\cite{shakson2025ssmcgm}.
\end{center}
\end{frame}

% ----------------------------------------------------
% Frame 8: Architecture Comparison (Visual First)
% ----------------------------------------------------
\begin{frame}{Architecture: TFT vs SSM-CGM}
    \vspace{0.5em}
    \begin{figure}
        \centering
        \begin{minipage}{0.48\textwidth}
            \centering
            \includegraphics[width=\linewidth]{TFT_architecture.png}
            \vspace{8pt}
            \textbf{Temporal Fusion Transformer (TFT)}
            Figure from Lim et al., Temporal Fusion Transformers Architecture (2021)~\cite{lim2021temporal}.
        \end{minipage}
        \hfill
        \begin{minipage}{0.48\textwidth}
            \centering
            \includegraphics[width=\linewidth]{SSM_architecture.png}
            \vspace{8pt}
            \textbf{SSM-CGM}         
            SSM-CGM Model Architecture from Shakson et al.~\cite{shakson2025ssmcgm}.
        \end{minipage}
        
        \vspace{0.1em}
    \end{figure}
    \vspace{0.8em}
\end{frame}

% ----------------------------------------------------
% Frame 9: Component Comparison Table
% ----------------------------------------------------
\begin{frame}{TFT $\to$ SSM-CGM: What Was Replaced and Why}

\begin{center}
\begin{tabular}{@{}p{3.2cm}@{\hspace{0.4cm}}p{2.6cm}@{\hspace{0.4cm}}p{3.6cm}@{\hspace{0.4cm}}p{3.6cm}@{}}
\toprule
\textbf{Component} & \textbf{Status} & \textbf{TFT (Lim et al., 2021)} & \textbf{SSM-CGM (Shakson et al., 2025)} \\
\midrule

Variable Selection (VSN)
& \textcolor{blue}{Kept}
& Selects relevant inputs
& Identical \\[6pt]

Gated Residual Blocks (GRBs)
& \textcolor{blue}{Kept}
& Gated residual mixing
& Identical \\[6pt]

Prediction Head
& \textcolor{blue}{Kept}
& Multi-horizon quantile head
& Identical \\[6pt]

\textbf{Local Temporal Processing}
& \textcolor{ForestGreen}{Replaced}
& LSTM encoder-decoder \newline
\scriptsize Sequential, slow, vanishing gradients
& \textbf{4$\times$ Mamba blocks} \newline
\scriptsize Parallel, linear-time, long-range memory \\[8pt]

\textbf{Global Context Modeling}
& \textcolor{ForestGreen}{Replaced}
& Multi-head attention \newline
\scriptsize $\mathcal{O}(L^2)$ complexity
& \textbf{Light Mamba layer} \newline
\scriptsize $\mathcal{O}(L)$ complexity \\

\bottomrule
\end{tabular}
\end{center}

\footnotesize
\textbf{Summary:} SSM-CGM swaps TFT's temporal backbone with Mamba \cite{gu2023mamba},
while preserving \textbf{all interpretability features} (VSN, quantile outputs).
\end{frame}

% ----------------------------------------------------
% Frame 10: TFT Formula Explanation (Now makes sense after seeing the problem)
% ----------------------------------------------------
\begin{frame}{TFT Prediction Formula — What Does It Really Mean?}

\begin{center}
\small
$\hat{y}_i(q, t, \tau) 
    = f_q\Big( \tau,\;
    \textcolor{ForestGreen}{y_{i,t-k:t}},\;
    \textcolor{violet}{z_{i,t-k:t}},\;
    \textcolor{BrickRed}{x_{i,t:t+\tau}},\;
    \textcolor{Orange}{s_i}
    \Big)$
\end{center}
 
\footnotesize

\begin{block}{}
\centering
\begin{tabular}{@{}lllp{7cm}@{}}
\toprule
\textbf{Symbol} & \textbf{Color} & \textbf{Input Type} & \textbf{Meaning + Link to TFT Architecture} \\
\midrule
$\hat{y}_i(q,t,\tau)$ 
    & — & Output 
    & Predicted glucose at quantile $q$ (e.g. 0.1/0.5/0.9), time $t$, horizon $\tau$ 
      $\to$ comes from the \alert{Multi-Horizon Quantile Head} \\[6pt]

$f_q(\cdot)$ 
    & — & Model 
    & The full TFT network (VSN + GRUs + LSTM Encoder/Decoder + Self-Attention) \\[6pt]

\textcolor{ForestGreen}{$y_{i,t-k:t}$} 
    & Green 
    & \textbf{Past observed target} 
    & Historical glucose values $\to$ processed by \alert{LSTM Encoder} and \alert{VSN (past)} \\[6pt]

\textcolor{violet}{$z_{i,t-k:t}$} 
    & Purple 
    & \textbf{Past known/unknown covariates} 
    & Past dynamic inputs (HR, steps, stress, etc.) $\to$ \alert{Variable Selection Network (past)} selects which ones matter \\[6pt]

\textcolor{BrickRed}{$x_{i,t:t+\tau}$} 
    & Red 
    & \textbf{Future known inputs} 
    & Time-of-day, day-of-week, predicted meal flags, holidays $\to$ fed directly to \alert{LSTM Decoder} and \alert{VSN (future)} \\[6pt]

\textcolor{Orange}{$s_i$} 
    & Yellow 
    & \textbf{Static covariates} 
    & Patient age, diabetes type, HbA1c, ID embedding $\to$ injected via \alert{Static Covariate Encoders} into every layer \\
\bottomrule
\end{tabular}
\end{block}
\end{frame}

\begin{frame}{TFT Prediction}
    
\normalsize
In plain English: \\
\alert{To predict your glucose in the next hour (with uncertainty), TFT intelligently combines:}\\
your recent glucose trend + past wearable data + known future events (meals, time of day) + who you are as a patient.
\end{frame}

% ----------------------------------------------------
% Frame 11: SSM-CGM Formula Explanation
% ----------------------------------------------------
\begin{frame}{SSM-CGM Prediction Formula — Simpler \& Faster than TFT}

\begin{center}
\small
$\hat{y}_{t+1:t+H} 
= f_\theta \Big( 
    \textcolor{ForestGreen}{y_{t-L+1:t}},\;
    \textcolor{violet}{x^{(hist)}_{t-L+1:t}},\;
    \textcolor{BrickRed}{x^{(fut)}_{t+1:t+H}},\;
    \textcolor{Orange}{s}
\Big)$
\end{center}

\footnotesize
\begin{block}{}
\centering
\begin{tabular}{@{}lllp{7cm}@{}}
\toprule
\textbf{Symbol} & \textbf{Color} & \textbf{Input Type} & \textbf{Meaning + Link to SSM-CGM Architecture} \\
\midrule

$\hat{y}_{t+1:t+H}$ 
    & — & Output 
    & Predicted glucose trajectory for next $H$ steps (e.g. next 60 min) 
      $\to$ output by the same \alert{Multi-Horizon Quantile Head} as TFT \\[6pt]

$f_\theta(\cdot)$ 
    & — & Model 
    & Full SSM-CGM network (= TFT's interpretable parts + \textbf{Mamba} backbone) \\[6pt]

\textcolor{ForestGreen}{$y_{t-L+1:t}$} 
    & Green 
    & Past glucose 
    & Last $L$ observed glucose values (e.g. past 48h) $\to$ processed by \alert{4 stacked Mamba blocks (local)} \\[6pt]

\textcolor{violet}{$x^{(hist)}_{t-L+1:t}$} 
    & Purple 
    & Past dynamic covariates 
    & Heart rate, respiration, steps, stress, meal flags, etc. $\to$ selected by \alert{VSN (past)} and fed into Mamba layers \\[6pt]

\textcolor{BrickRed}{$x^{(fut)}_{t+1:t+H}$} 
    & Red 
    & Future known inputs 
    & Time-of-day, predicted future meal flags, circadian features $\to$ selected by \alert{VSN (future)} and processed by \alert{Light Mamba layer (global)} \\[6pt]

\textcolor{Orange}{$s$} 
    & Yellow 
    & Static patient info 
    & Age, diabetes type, HbA1c, patient ID embedding $\to$ injected via \alert{static covariate encoders} into every layer (unchanged from TFT) \\
\bottomrule
\end{tabular}
\end{block}

\vspace{0.8em}
\normalsize
\textbf{In plain English:} \\
SSM-CGM uses \alert{exactly the same inputs and interpretability tools as TFT} (VSN, static encoders, quantile head) — \\
but replaces the slow LSTM + quadratic self-attention with \textbf{linear-time Mamba blocks}.

Result: same clinical interpretability, better long-range modeling, 6–7\% lower error on 24h–48h context.

\vspace{0.8em}
\footnotesize
Formula and architecture from Shakson et al.~\cite{shakson2025ssmcgm}.
\end{frame}

% ----------------------------------------------------
% Frame 12: Understanding SSMs (Math)
% ----------------------------------------------------
\begin{frame}{Understanding State-Space Models: The Math Made Simple}
    \textbf{Intuition:} A State-Space Model maintains a "hidden memory" ($\mathbf{h}_t$) that gets updated as new data ($\mathbf{x}_t$) arrives, then produces predictions ($\mathbf{z}_t$).
    
    \vspace{1em}
    \begin{columns}
        \begin{column}{0.5\textwidth}
            \begin{block}{Step 1: Update Memory}
                $$\mathbf{h}_{t+1} = \mathbf{A}_t \mathbf{h}_t + \mathbf{B}_t \mathbf{x}_t$$
                \begin{itemize}
                    \item $\mathbf{h}_{t+1}$: New memory state
                    \item $\mathbf{A}_t$: How much of the old state to keep ("remember")
                    \item $\mathbf{B}_t$: How much new input to incorporate
                    \item $\mathbf{x}_t$: Current input (e.g., current glucose, heart rate)
                \end{itemize}
            \end{block}
        \end{column}
        
        \begin{column}{0.5\textwidth}
            \begin{block}{Step 2: Generate Prediction}
                $$\mathbf{z}_t = \mathbf{C}_t \mathbf{h}_t$$
                \begin{itemize}
                    \item $\mathbf{z}_t$: Model output (predicted glucose trajectory)
                    \item $\mathbf{C}_t$: Readout matrix (extracts relevant info from memory)
                    \item $\mathbf{h}_t$: Current memory state
                \end{itemize}
            \end{block}
        \end{column}
    \end{columns}
\end{frame}

\begin{frame}
    \vspace{1em}
    \begin{alertblock}{Mamba's Breakthrough: Selectivity}
        In traditional SSMs, $\mathbf{A}$, $\mathbf{B}$, and $\mathbf{C}$ are fixed. In \textbf{Mamba}, these matrices are \textbf{input-dependent} and change dynamically:
        \begin{itemize}
            \item The model \textbf{learns} what to remember or forget based on the current context
            \item Enables personalized tracking: different patterns for different individuals
            \item Example: After a meal flag appears, $\mathbf{A}_t$ might "remember" this longer; during stable overnight periods, it might "forget" faster
        \end{itemize}
    \end{alertblock}
\end{frame}

% ----------------------------------------------------
% TRANSITION 4: From Architecture to Results
% ----------------------------------------------------
\begin{frame}[plain]
    \centering
    \vfill
    \begin{beamercolorbox}[center,shadow=true,rounded=true]{title}
        \usebeamerfont{title}
        \textbf{Part IV: Does It Work?}\\
        \vspace{0.5em}
        \large Experimental validation and performance analysis
    \end{beamercolorbox}
    \vfill
    
    \vspace{1em}
    \normalsize
    \begin{block}{What We'll Evaluate}
        \begin{itemize}
            \item \textbf{Accuracy:} Does SSM-CGM outperform TFT?
            \item \textbf{Interpretability:} Can we understand \textit{what} and \textit{when} features matter?
            \item \textbf{Actionability:} Can we simulate "what-if" scenarios?
        \end{itemize}
    \end{block}
    
    \vspace{1em}
    \textit{Let's examine the training setup and results...}
\end{frame}

% ----------------------------------------------------
% Frame 13: Training Setup
% ----------------------------------------------------
\begin{frame}{Training \& Experimental Setup}
    \begin{block}{Forecasting Task}
        \textbf{Objective:} Predict the next \textbf{1 hour} (12 time steps at 5-min resolution) of glucose trajectory \\
        \textbf{Input Context:} Used 12h, 24h, or 48h of historical data
        \begin{itemize}
            \item 12h = 144 time steps; 24h = 288 steps; 48h = 576 steps
        \end{itemize}
        \textbf{Data Split:} Last 2 hours withheld (penultimate hour for validation, final hour for testing)
    \end{block}
    
    \begin{columns}
        \begin{column}{0.5\textwidth}
            \begin{block}{Evaluation Metrics}
                \textbf{Mean Absolute Error (MAE):}
                $$\text{MAE} = \frac{1}{N} \sum_{i=1}^{N} |y_i - \hat{y}_i|$$
                \begin{itemize}
                    \item Average magnitude of errors (mg/dL)
                    \item Intuitive: "on average, predictions are off by X mg/dL"
                    \item Treats all errors equally
                \end{itemize}
            \end{block}
        \end{column}
        
        \begin{column}{0.5\textwidth}
            \begin{block}{}
                \textbf{Root Mean Squared Error (RMSE):}
                $$\text{RMSE} = \sqrt{\frac{1}{N} \sum_{i=1}^{N} (y_i - \hat{y}_i)^2}$$
                \begin{itemize}
                    \item \textbf{Penalizes large errors} more heavily
                    \item Critical for clinical safety: large errors (e.g., missing hypoglycemia) are dangerous
                \end{itemize}
            \end{block}
        \end{column}
    \end{columns}
    
    \vspace{0.05em}
    \begin{block}{Primary Training Loss: Quantile Loss}
        Trained with quantile regression ($q \in \{0.1, 0.5, 0.9\}$) to output \textbf{prediction intervals} rather than single-point estimates:
        \begin{itemize}
            \item $q=0.5$ (median): central prediction
            \item $q=0.1, 0.9$: lower/upper bounds capturing uncertainty
            \item \textbf{Clinical Value:} Enables risk assessment -- wider intervals signal higher uncertainty
        \end{itemize}
    \end{block}
\end{frame}

% ----------------------------------------------------
% Frame 14: Results
% ----------------------------------------------------
\begin{frame}{Results: SSM-CGM Outperforms TFT}
    \begin{block}{Experimental Fairness}
        \begin{itemize}
            \item Identical preprocessing, features, and data splits
            \item Matched parameter counts: SSM-CGM (1.76M) vs TFT (1.79M)
            \item Same hyperparameter search space and training protocol
            \item Both models trained on 4$\times$A100 GPUs with identical optimization
        \end{itemize}
    \end{block}
\end{frame}

\begin{frame}
    \begin{table}
        \centering
        \caption{Performance comparison for 1-hour glucose forecasting (mean with 95\% CI)}
        \begin{tabular}{lccccc}
            \toprule
            \textbf{Context} & \textbf{Model} & \textbf{MAE (mg/dL)} & \textbf{RMSE (mg/dL)} & \textbf{$\Delta$ MAE} \\
            \midrule
            12h & TFT & 6.07 (5.76, 6.37) & 7.32 (6.96, 7.69) & -- \\
                & SSM-CGM & 5.97 (5.68, 6.26) & 7.17 (6.83, 7.52) & \textcolor{ForestGreen}{-1.6\%} \\
            \midrule
            24h & TFT & 6.06 (5.77, 6.36) & 7.28 (6.93, 7.63) & -- \\
                & SSM-CGM & 5.66 (5.39, 5.93) & 6.86 (6.53, 7.19) & \textcolor{ForestGreen}{-6.6\%} \\
            \midrule
            48h & TFT & 5.77 (5.50, 6.03) & 6.92 (6.61, 7.24) & -- \\
                & SSM-CGM & 5.40 (5.15, 5.65) & 6.56 (6.26, 6.87) & \textcolor{ForestGreen}{-6.5\%} \\
            \bottomrule
        \end{tabular}
    \end{table}
\end{frame}

\begin{frame}{Key Observations from Performance Comparison}
      
     \begin{block}{Key Observations}
        \begin{itemize}
            \item SSM-CGM consistently outperforms TFT across all context lengths
            \item \textbf{Performance gap widens with longer context} (48h: 6.5\% improvement vs 12h: 1.6\%)
            \item Demonstrates SSM-CGM's superior ability to leverage long-term physiological history
            \item Within SSM-CGM: $\sim$5\% improvement with each doubling of context (12h$\rightarrow$24h$\rightarrow$48h)
        \end{itemize}
    \end{block}
\end{frame}

% ----------------------------------------------------
% Frame 15: Interpretability
% ----------------------------------------------------
\begin{frame}{Interpretability: Opening the Black Box for Clinical Use}
    \textbf{Why Interpretability Matters:} Clinicians need to understand \textbf{\textit{why}} a model makes predictions to trust and act on them.
   
  
    \begin{columns}[T]   % ← the [T] is the key: top-align both columns
        \begin{column}{0.49\textwidth}
            \begin{block}{1. Variable Importance (What Matters?)}
                \textbf{Method:} Variable Selection Networks (VSN) assign interpretable weights to each input feature
               
                \vspace{0.5em}
                \textbf{Key Findings:}
                \begin{itemize}
                    \item \textbf{Top predictors:} Past CGM values, predicted meal flags, heart rate, respiration rate
                    \item Reflects known physiology: meals cause glucose spikes, activity modulates glucose uptake
                   
                \end{itemize}
            \end{block}
        \end{column}
        \hfill
        \begin{column}{0.49\textwidth}
            \begin{block}{2. Temporal Importance (When Matters?)}
                \textbf{Method:} ``Hidden attention'' from SSM dynamics reveals which past time steps influence each prediction
               
                \vspace{0.5em}   % ← this line was missing → caused the misalignment
                \textbf{Key Findings:}
                \begin{itemize}
                    \item \textbf{Cohort-level pattern:} Recent history dominates (expected for smooth CGM)
                    \item \textbf{Individual-specific peaks:} Unique lag patterns per person — model learns personalized memory dependencies
                 
                \end{itemize}
            \end{block}
        \end{column}
    \end{columns}
    
    \vspace{0.8em}
    \begin{alertblock}{Clinical Translation}
        These two complementary interpretability signals allow clinicians to answer both \textbf{``what''} matters (features) and \textbf{``when''} it mattered (time lags) → essential for trust and real-world adoption.
    \end{alertblock}
\end{frame}

\begin{frame}
    
    \begin{columns}[T]   % ← the [T] is the key: top-align both columns
        \begin{column}{0.49\textwidth}
            \begin{block}{1. Variable Importance (What Matters?)}
              
                \vspace{0.5em}
                \textbf{Key Findings:}
                \begin{itemize}
                
                    \item Wearable features (HR, RR, stress) capture circadian and behavioral patterns
                    \item \textbf{Static features:} Participant ID most important (individual variation), followed by diabetes status and glucose variability
                \end{itemize}
            \end{block}
        \end{column}
        \hfill
        \begin{column}{0.49\textwidth}
            \begin{block}{2. Temporal Importance (When Matters?)}
                % ← this line was missing → caused the misalignment
                \textbf{Key Findings:}
                \begin{itemize}
                    
                    \item \textbf{Layered specialization:} Early Mamba layers focus on recent lags (short-term), deeper layers capture distant events (circadian)
                \end{itemize}
            \end{block}
        \end{column}
    \end{columns}
     \begin{alertblock}{Clinical Translation}
        These interpretability features enable clinicians to:
        \begin{itemize}
            \item Verify model reasoning aligns with physiological knowledge
            \item Identify which patient-specific factors drive predictions
            \item Build trust in model recommendations
        \end{itemize}
    \end{alertblock}
  
\end{frame}

\begin{frame}{Variable Importance (VSN) Evolves Over Time}

\begin{figure}
    \centering
    \begin{minipage}{0.495\textwidth}
        \centering
        \includegraphics[width=\linewidth]{encoder_heatmap_vars_importance.png}
         \textbf{Encoder VSN weights} \\
        \small (past → present context)
 
    \end{minipage}
    \hfill
    \begin{minipage}{0.495\textwidth}
        \centering
        \includegraphics[width=\linewidth]{heatmap_vars_importance.png}
         \textbf{Decoder VSN weights} \\
        \small (future prediction horizons)
     \end{minipage}
    \begin{center}
    \tiny Figures adapted from Shakson et al.~\cite{shakson2025ssmcgm} (Figure 9: VSN variable by time importance for the encoder and the decoder)

    \end{center}
     \caption*{\textbf{Dynamic Variable Selection Network (VSN) importance heatmaps} \\
             Each row = one input feature (CGM, meal flag, HR, RR, stress, static covariates, etc.). \\
             Color intensity shows how much the model relies on that feature at each time step. \\
             The model automatically learns that, e.g., meal flags matter most right after eating, \\
             heart rate matters more during activity, and static patient features remain influential throughout.}
\end{figure}

\vspace{1em}
\footnotesize
.

\end{frame}

% ----------------------------------------------------
% Frame 16: Counterfactual Forecasting
% ----------------------------------------------------
\begin{frame}{Counterfactual Forecasting: Answering "What-If?" Questions}
    \begin{block}{The Innovation}
        SSM-CGM enables \textbf{counterfactual reasoning}: simulate how planned changes in behavior (via wearable signals) would affect future glucose.
    \end{block}
    
    \begin{block}{Methodology}
        \textbf{Theoretical Foundation:} Sequential g-formula under standard causal assumptions:
        \begin{itemize}
            \item Consistency: Planned actions match observed measurements
            \item Sequential ignorability: No unmeasured confounders (given rich history)
            \item Positivity: All planned actions are physiologically plausible
        \end{itemize}
        
        \vspace{0.5em}
        \textbf{Implementation:} Intervene on future covariates (HR, RR) in model input, holding all else equal:
        \begin{center}
            \textit{"What if my heart rate increases by +2 SD over the next hour?"}
        \end{center}
    \end{block}
\end{frame}  

\begin{frame}
    \begin{columns}
        \begin{column}{0.5\textwidth}
            \begin{block}{Simulated Effects}
                \textbf{Heart Rate (HR) +2 SD:}
                \begin{itemize}
                    \item Effect: \textcolor{ForestGreen}{$-5.0$ mg/dL} glucose drop
                    \item Mechanism: Simulates physical activity (daytime)
                    \item Correlated with cardiovascular fitness markers
                \end{itemize}
                
                \vspace{0.5em}
                \textbf{Respiration Rate (RR) +2 SD:}
                \begin{itemize}
                    \item Effect: \textcolor{BrickRed}{$+1.2$ mg/dL} glucose increase
                    \item Mechanism: Simulates impaired sleep breathing (nighttime)
                    \item Correlated with adiposity and inflammation
                \end{itemize}
            \end{block}
        \end{column}
        
        \begin{column}{0.5\textwidth}
            \begin{block}{Clinical Implications}
                \begin{itemize}
                    \item \textbf{Personalized Guidance:} Model translates physiological signals into glucose outcomes
                    \item \textbf{Proactive Planning:} Patients can see predicted impact of planned activities
                    \item \textbf{Example Use Cases:}
                    \begin{itemize}
                        \item "Will a 30-min walk prevent a post-meal spike?"
                        \item "How will poor sleep quality affect tomorrow's glucose?"
                    \end{itemize}
                    \item \textbf{Validation:} Effects align with known metabolic physiology
                \end{itemize}
            \end{block}
        \end{column}
    \end{columns}
    
\end{frame}

\begin{frame}{Counterfactual Forecasting: ``What-If'' Scenarios}

\begin{figure}
    \centering
    \includegraphics[width=0.96\linewidth]{glucose_counterfactual.png}
    \tiny Figure adapted from Shakson et al.~\cite{shakson2025ssmcgm} 

    \vspace{0.8em}
    \caption*{\textbf{Counterfactual glucose trajectories under +2 SD interventions} \\
             \textcolor{ForestGreen}{↑ Heart rate (+2 SD)} → ∼5 mg/dL reduction \quad 
             \textcolor{BrickRed}{↑ Respiration rate (+2 SD)} → ∼1–2 mg/dL increase \\
             Interventions applied at t = 0 (vertical dashed line). Shaded = 90\% PI.}
\end{figure}

\vspace{1.2em}
\footnotesize
Figure adapted from Shakson et al.~\cite{shakson2025ssmcgm}.

\end{frame}

% ----------------------------------------------------
% Frame 17: Conclusion
% ----------------------------------------------------
\begin{frame}{Conclusion: Advancing Personalized Diabetes Care}
    \begin{block}{Key Achievements}
        \begin{enumerate}
            \item \textbf{State-of-the-Art Accuracy:}
            \begin{itemize}
                \item SSM-CGM outperforms TFT, especially for long-context forecasting (6.5\% MAE improvement at 48h)
                \item Leverages Mamba's linear-time complexity for efficient long-sequence modeling
            \end{itemize}
            
            \item \textbf{Clinical Interpretability:}
            \begin{itemize}
                \item Variable importance: identifies which signals drive predictions
                \item Temporal attribution: reveals when past events influence forecasts
                \item Personalized insights: captures individual-specific patterns
            \end{itemize}
            
            \item \textbf{Actionable Counterfactuals:}
            \begin{itemize}
                \item Enables "what-if" scenario planning
                \item Links physiological inputs (HR, RR) to glucose outcomes
                \item Supports proactive behavioral interventions
            \end{itemize}
        \end{enumerate}
    \end{block}
\end{frame}

\begin{frame}
    \begin{block}{Broader Impact}
        \begin{itemize}
            \item \textbf{Computational Efficiency:} Linear-time complexity enables real-time, on-device deployment
            \item \textbf{Trust \& Adoption:} Interpretability builds clinician confidence in AI-driven recommendations
            \item \textbf{Paradigm Shift:} From reactive glucose monitoring to \textbf{predictive, personalized} diabetes management
        \end{itemize}
    \end{block}
    
    \vspace{0.5em}
    \begin{center}
        \textit{SSM-CGM represents a critical step toward integrating cutting-edge AI with the real-world needs of personalized healthcare.}
    \end{center}
\end{frame}

% ----------------------------------------------------
% Frame 18: Future Work
% ----------------------------------------------------
\begin{frame}{Future Directions}
    \begin{block}{1. Data Integration \& Refinement}
        \begin{itemize}
            \item Ground-truth meal data (timing, size, macros)
            \item Medication tracking (insulin, metformin, GLP-1)
            \item Extended horizons (2--4 hours)
            \item Multivariate forecasting (glucose + wearables)
        \end{itemize}
    \end{block}
    
    \begin{block}{2. Deployment \& Validation}
        \begin{itemize}
            \item On-device implementation for real-time alerts
            \item Streaming processing for continuous updates
            \item Clinical trials for TIR impact validation
            \item Multi-cohort generalization testing
        \end{itemize}
    \end{block}
\end{frame}

\begin{frame}
    \begin{block}{3. From Prediction to Action}
        \begin{itemize}
            \item Reinforcement learning for optimal policies
            \item Multivariate interventions (meal + exercise + sleep)
            \item Adaptive recommendations ("Take 1U insulin" / "Walk 15 min")
            \item Integration with automated insulin delivery
        \end{itemize}
    \end{block}
\end{frame}

\begin{frame}[allowframebreaks]{References}
\frametitle{References}
\begin{thebibliography}{10}
\setlength{\itemsep}{0.5em}

\bibitem{shakson2025ssmcgm}
Shakson Isaac, Yentl Collin, and Chirag Patel.
\newblock {\em SSM-CGM: Interpretable State-Space Forecasting Model of Continuous Glucose Monitoring for Personalized Diabetes Management}.
\newblock Preprint arXiv:2502.12345, February 2025.
\newblock \textcolor{blue}{\url{https://arxiv.org/abs/2502.12345}}

\bibitem{lim2021temporal}
Bryan Lim, Sercan Ö. Arık, Nicolas Loeff, and Tomas Pfister.
\newblock Temporal Fusion Transformers for Interpretable Multi-Horizon Time Series Forecasting.
\newblock \emph{International Journal of Forecasting}, 37(4):1748--1764, 2021.

\bibitem{gu2023mamba}
Albert Gu and Tri Dao.
\newblock Mamba: Linear-Time Sequence Modeling with Selective State Spaces.
\newblock \emph{International Conference on Learning Representations (ICLR)}, 2023.

\end{thebibliography}
\end{frame}

% ----------------------------------------------------
% Frame 19: Q&A
% ----------------------------------------------------
\begin{frame}[standout]
    \centering
    \Huge Thank You! \\
    \vspace{1cm}
    \Large Questions \& Discussion
    
    \vspace{2cm}
    \normalsize
\end{frame}

\end{document}